\scanpage{38}
$Z_i$ of $S$ and the action of $G_i$ on $E$.\inserted{When $S$ is again arbitrary,}This shows in a \unsure{rather|rahher} concrete way
what it means, \inserted{locally(for étale topology)} for a covering $S'$ of $S$, to be associated to some $M$: it
means that locally it is describable by a representation of an inertia\deleted{l}
group on a finite set, in the obvious way; and the extra structure
on one $S'$ involved by deducing it from some $S$, can be described by
saying that we give such local representations of $S'$, in such a way
that they match together i.e. satisfy some rather evident compatibility
condition (This description makes a sense whenever the functor ``\unsure{geometri}
realization'' is faithful - not necessarily fully faithful). The analogous
description holds for Galois coverings with Galois group $\unsure{G|H}$, namely
in the strictly local case, for a fixed $Z_i$ , the\deleted{y are classified by
systems (} category of these is equivalent to the category of finite
sets $E$ on which $G_i$ operates on the left, $H$ on the right, $G_i$ and $H$ commuting,
in such a way that $E$ is a principal homogeneous space (or, as we
say, a torsor) under $H$; up to isomorphism, they are classified by
$\operatorname{Hom}(G_i,H)/H$ , group homomorphisms mod. interior automorphism. Using \deleted{th}
this description, one easily proves the criterion I indicated for the
``geometric\deleted{\unsure{-|r}}\unsure{reallization|realization} -functor'' to be fully faithful on the category
of \c{\u{G}}alois coverings of type $\underline{R}$ ; indeed, as usual the question is a \unsure{local|pocal}
one for the étale topology, which allows one to \deleted{restr} reduce to\unsure{'} the
stricly local case.

Your interpretation of the Kummer case in the final \unsure{f|s}\unsure{o|i}rmulation
of example 3 is indeed the one I had in \unsure{mihd|mind}. Also, when I wrote $n'\quad n$,
I meant of course the order relation of divisibility (it may be convevient
to introduce this order relation explicitly, for simplicity of
notations).
% archive sheet number: 37
